\section{Dzienniczek pacjenta / ePRO}

% SLIDE ID: sec04-goal
\BlockGoalFrame{\SlideID{sec04-goal}}{
  \item pokazać, czym ePRO różni się od klasycznego zbierania danych w ośrodku,
  \item omówić ryzyka użyteczności i kompletności danych,
  \item połączyć perspektywę pacjenta z wymaganiami badania.
}

% SLIDE ID: sec04-core
\begin{frame}{ePRO jako źródło danych od pacjenta}
  \SlideID{sec04-core}
  \begin{itemize}
    \item ePRO umożliwia zbieranie danych bezpośrednio od uczestnika badania.
    \item Rozwiązania te zwiększają częstotliwość i aktualność danych.
    \item Kluczowe są użyteczność interfejsu, zgodność i kontrola kompletności wpisów.
  \end{itemize}
\end{frame}

% SLIDE ID: sec04-case
\CaseStudyFrame{Brak wpisów w dzienniczku pacjenta}{
  W badaniu hybrydowym część pacjentów regularnie opuszcza wpisy wieczorne mimo aktywnych przypomnień.
}{
  Czy problem leży w projekcie interfejsu, harmonogramie wpisów, komunikacji z pacjentem czy logice monitoringu?
}
